% ==============================================================
% Section 1: Introduction
% Author: Quân (RW)
% File:   paper/sections/01_intro.tex
% ==============================================================

\section{Introduction}\label{sec:intro}

Bug reports are the primary channel through which defects are
communicated in open-source software (OSS) projects.  A well-structured
report allows developers to understand what went wrong, compare it
against expected behavior, and reproduce the defect reliably.  Prior
research has established that the three components most critical to
this process are the \emph{Observed Behavior} (OB), the \emph{Expected
Behavior} (EB), and the \emph{Steps to Reproduce}
(S2R)~\cite{song2020bee, mahmud2025astrobr, akyol2026improbr}.  When
any of these is absent or ambiguous, developer comprehension suffers
and resolution time increases.

Despite this, many real-world bug reports remain
incomplete~\cite{fan2020chaff}.  Reviewing each incoming report
manually at scale is expensive and does not fit the workflows of most
OSS projects.  This motivates automated tools that can screen reports
before they reach the development queue.

Large language models (LLMs) have been applied to bug-report tasks,
including quality classification and improvement~\cite{bo2024chatbr,
acharya2025enhance}.  However, as reviewed in
Section~\ref{sec:related}, \textbf{no existing study measures Cohen's
Kappa between an LLM and human annotators for the combined
OB\,+\,EB\,+\,S2R quality dimensions on general GitHub repositories}.
Studies that employ inter-rater agreement metrics are either restricted
to platform-specific contexts (Android, Mojira) or measure only
human--human agreement during ground-truth construction.

This study addresses that gap.  We evaluate \emph{GPT-4o mini} as an
automated quality assessor for GitHub bug reports, measuring its label
agreement with an annotator consensus using Cohen's Kappa ($\kappa$)
across OB, EB, and S2R.  GPT-4o mini was selected because a quality
gate must run on every incoming issue: at approximately
US\$0.0003 per report (Section~\ref{sec:method:repro}), a lightweight
model is the realistic deployment candidate.  The evaluation covers 30
annotated issues from three OSS repositories (pandas, scikit-learn,
VS~Code).

The research question is:

\begin{quote}
\textit{How reliably can GPT-4o mini automatically evaluate bug report
quality (OB, EB, and S2R sufficiency) compared to annotator consensus,
as measured by Cohen's Kappa, on a dataset of open-source GitHub bug
reports?}
\end{quote}

The main contributions of this paper are:

\begin{itemize}
  \item A structured annotation rubric with four sufficiency labels
        (\textit{Sufficient}, \textit{Ambiguous}, \textit{Missing},
        \textit{Incorrect}) for OB, EB, and S2R in GitHub bug reports;
  \item An empirical evaluation of GPT-4o mini's agreement with
        annotator consensus across all three dimensions, with Cohen's
        Kappa as the primary metric;
  \item A dimension-level analysis of the model's strengths and
        limitations, with practical guidance on integrating LLM-based
        screening into bug-report triage pipelines.
\end{itemize}

The remainder of this paper is organized as follows.
Section~\ref{sec:related} reviews related work.
Section~\ref{sec:method} describes the methodology.
Section~\ref{sec:results} presents the results.
Section~\ref{sec:discussion} discusses the findings.
Section~\ref{sec:threats} addresses threats to validity, and
Section~\ref{sec:conclusion} concludes the paper.

